% uid: ""
% school: "lycee"
% level: "3eme"
% section: ["Economie et Gestion"]
% chapter: "Statistiques"
% title: "Équilibre du marché : offre et demande"
% difficulty: "Difficile"
% points: null
% tags: ["statistiques", "serie-double", "droite-de-mayer", "offre-demande", "equilibre", "modelisation"]
% provenance: "own"
% dressing: "deux ajustements"
% rung: 4
% batch: "2026-11-20-stat-3eco-A"

% ---STATEMENT---
Pour un produit agricole, le prix $x$ est exprimé en dinars le kilo. Des enquêtes ont fourni :
\begin{itemize}
  \item côté \textbf{demande} : un ajustement affine donne $y_d=-6{,}5\,x+109{,}5$ (en tonnes) ;
  \item côté \textbf{offre} : les producteurs ont proposé les quantités $y_s$ (en tonnes) suivantes :
\end{itemize}
\begin{center}
\begin{tabular}{|l|c|c|c|c|c|c|c|c|}
\hline
Prix $x_i$ & $2$ & $3$ & $4$ & $5$ & $6$ & $7$ & $8$ & $9$ \\ \hline
Offre $y_{s,i}$ & $30$ & $38$ & $44$ & $53$ & $58$ & $64$ & $70$ & $77$ \\ \hline
\end{tabular}
\end{center}

\begin{enumerate}
  \item Par la méthode de Mayer (deux sous-nuages de $4$ points), déterminer une équation $y_s=ax+b$ de la droite d'ajustement de l'offre.
  \item Résoudre $y_s=y_d$. En déduire le prix d'équilibre et la quantité échangée à l'équilibre.
  \item Pour un prix de $5$ dinars, calculer $y_s$ et $y_d$. Y a-t-il pénurie ou excédent ? De combien de tonnes ?
  \item Interpréter graphiquement le prix d'équilibre.
\end{enumerate}

% ---SOLUTION---
\begin{enumerate}
  \item Sous-nuages : prix $2$ à $5$, puis prix $6$ à $9$.
  $$H_1:\ \overline{x}=3{,}5,\ \overline{y}=\frac{30+38+44+53}4=\frac{165}4=41{,}25\qquad
  H_2:\ \overline{x}=7{,}5,\ \overline{y}=\frac{58+64+70+77}4=\frac{269}4=67{,}25.$$
  Pente : $a=\dfrac{67{,}25-41{,}25}{7{,}5-3{,}5}=\dfrac{26}4=6{,}5$. Puis $41{,}25=6{,}5\times3{,}5+b$ donne $b=41{,}25-22{,}75=18{,}5$.
  $$ \ib{$y_s=6{,}5\,x+18{,}5$}.$$
  \item L'équilibre est atteint quand l'offre est égale à la demande :
  $$6{,}5\,x+18{,}5=-6{,}5\,x+109{,}5\iff13\,x=91\iff x=7.$$
  Quantité échangée : $y_s(7)=6{,}5\times7+18{,}5=45{,}5+18{,}5=64$. \emph{Vérification} avec la demande : $-6{,}5\times7+109{,}5=64$ $\checkmark$

  Prix d'équilibre : $\ib{$7$ dinars le kilo}$ ; quantité d'équilibre : $\ib{$64$ tonnes}$.
  \item Pour $x=5$ :
  $$y_s=6{,}5\times5+18{,}5=51\qquad y_d=-6{,}5\times5+109{,}5=77.$$
  La demande ($77$ t) dépasse l'offre ($51$ t) : \textbf{pénurie} de $77-51=26$ tonnes. Les acheteurs sont prêts à payer plus cher : le prix monte vers $7$ dinars.
  \item Dans un repère $(x\,;y)$, la droite d'offre est croissante, celle de demande décroissante. Elles se coupent au point $(7\,;64)$ : son abscisse est le prix d'équilibre et son ordonnée la quantité échangée.
\end{enumerate}

\textbf{Idée :} le prix d'équilibre est l'abscisse du point d'intersection des deux droites d'ajustement.
